\documentclass[11pt]{article}

% Change "review" to "final" to generate the final (sometimes called camera-ready) version.
% Change to "preprint" to generate a non-anonymous version with page numbers.
\usepackage[review]{acl}

% Standard package includes
\usepackage{times}
\usepackage{latexsym}

% For proper rendering and hyphenation of words containing Latin characters (including in bib files)
\usepackage[T1]{fontenc}
% For Vietnamese characters
% \usepackage[T5]{fontenc}
% See https://www.latex-project.org/help/documentation/encguide.pdf for other character sets

% This assumes your files are encoded as UTF8
\usepackage[utf8]{inputenc}

% This is not strictly necessary, and may be commented out,
% but it will improve the layout of the manuscript,
% and will typically save some space.
\usepackage{microtype}

% This is also not strictly necessary, and may be commented out.
% However, it will improve the aesthetics of text in
% the typewriter font.
\usepackage{inconsolata}

%Including images in your LaTeX document requires adding
%additional package(s)
\usepackage{graphicx}

% If the title and author information does not fit in the area allocated, uncomment the following
%
%\setlength\titlebox{<dim>}
%
% and set <dim> to something 5cm or larger.

\title{LocAgent: A Semi-Agentic System For Software Issue Localization}

% Author information can be set in various styles:
% For several authors from the same institution:
% \author{Author 1 \and ... \and Author n \\
%         Address line \\ ... \\ Address line}
% if the names do not fit well on one line use
%         Author 1 \\ {\bf Author 2} \\ ... \\ {\bf Author n} \\
% For authors from different institutions:
% \author{Author 1 \\ Address line \\  ... \\ Address line
%         \And  ... \And
%         Author n \\ Address line \\ ... \\ Address line}
% To start a separate ``row'' of authors use \AND, as in
% \author{Author 1 \\ Address line \\  ... \\ Address line
%         \AND
%         Author 2 \\ Address line \\ ... \\ Address line \And
%         Author 3 \\ Address line \\ ... \\ Address line}

\author{First Author \\
  Affiliation / Address line 1 \\
  Affiliation / Address line 2 \\
  Affiliation / Address line 3 \\
  \texttt{email@domain} \\\And
  Second Author \\
  Affiliation / Address line 1 \\
  Affiliation / Address line 2 \\
  Affiliation / Address line 3 \\
  \texttt{email@domain} \\}

%\author{
%  \textbf{First Author\textsuperscript{1}},
%  \textbf{Second Author\textsuperscript{1,2}},
%  \textbf{Third T. Author\textsuperscript{1}},
%  \textbf{Fourth Author\textsuperscript{1}},
%\\
%  \textbf{Fifth Author\textsuperscript{1,2}},
%  \textbf{Sixth Author\textsuperscript{1}},
%  \textbf{Seventh Author\textsuperscript{1}},
%  \textbf{Eighth Author \textsuperscript{1,2,3,4}},
%\\
%  \textbf{Ninth Author\textsuperscript{1}},
%  \textbf{Tenth Author\textsuperscript{1}},
%  \textbf{Eleventh E. Author\textsuperscript{1,2,3,4,5}},
%  \textbf{Twelfth Author\textsuperscript{1}},
%\\
%  \textbf{Thirteenth Author\textsuperscript{3}},
%  \textbf{Fourteenth F. Author\textsuperscript{2,4}},
%  \textbf{Fifteenth Author\textsuperscript{1}},
%  \textbf{Sixteenth Author\textsuperscript{1}},
%\\
%  \textbf{Seventeenth S. Author\textsuperscript{4,5}},
%  \textbf{Eighteenth Author\textsuperscript{3,4}},
%  \textbf{Nineteenth N. Author\textsuperscript{2,5}},
%  \textbf{Twentieth Author\textsuperscript{1}}
%\\
%\\
%  \textsuperscript{1}Affiliation 1,
%  \textsuperscript{2}Affiliation 2,
%  \textsuperscript{3}Affiliation 3,
%  \textsuperscript{4}Affiliation 4,
%  \textsuperscript{5}Affiliation 5
%\\
%  \small{
%    \textbf{Correspondence:} \href{mailto:email@domain}{email@domain}
%  }
%}

\begin{document}
\maketitle
\begin{abstract}
pass
\end{abstract}


\section{Introduction}
Repository-level code generation tackles the task of, given an issue or feature description and a codebase/repository, proposing a patch in the correct location to fix the issue or implement the feature.
Existing systems \citep{xia2024agentlessdemystifyingllmbasedsoftware, wang2025openhandsopenplatformai, yang2024sweagentagentcomputerinterfacesenable} typically follow a 3-stage pipeline: a \textbf{localizer} to identify the affected code, a \textbf{patch generator} to propose a fix, and a \textbf{verifier} to check correctness.
Trivially the performance of each of these components is bounded by the performance of the previous component, i.e. for patch generator to be successfully the localizer needs to provide the correct spots of the edit.
Agentless \citep{jimenez2024swebenchlanguagemodelsresolve} achieves ~70\% file-level localization ( lower function and line-level) correctness using GPT-4o on SWE-Bench Lite \citep{jimenez2024swebench} nd 32\% overall resolve rate.
To maximize the localization systems performance we focus on the localization component in the scope of this work, later on evaluating our improved localization system in an end to end setup to show the effectiveness.

Issue localization—the task of pinpointing the code that must change to fix a software problem is essential yet notoriously difficult.
The gap between natural-language descriptions and the underlying issue requires multi-hop reasoning over code and its dependencies.
Current LLM-based agents address this by augmenting themselves with repository retrieval and inspection tools.
In general localization systems are expected to provide fine grained locations of edit e.g. file-path(s) in the repository and the number of lines of the chunk(s).
Previous works vary on the level of granularity of the provided locations, such as SweRank \citep{reddy2025sweranksoftwareissuelocalization} performs function-level localization, nevertheless not all patches are within function changes, e.g. a change in the imports is required.
There are essentially 2 core approaches to issue localization.
In this work we also try to answer to the question of wether we need to pin-point the exact line numbers or providing the files is already enough.
Moreover providing a narrow chunk/window in a file might not provide enough context to the patch generator model about the file/class dependencies.

\textbf{Retrieval Based} methods – using encoder-only models to extract a latent representation of code chunks in the repo and match them with the issue description based on vector similarity metrics.
This method is limited due to the earlier mentioned gap between natural-language description and code-chunks.
Moreover it multiplicativally increases the overall cost of repository-level codegen tasks (e.g. SWE-Bench \citet{jimenez2024swebenchlanguagemodelsresolve}), since the output of the localization step should be fed to the patch generator and the output of the later to the verifier module.

\textbf{LLM-based} – where an agentic model is mean to perform various tool calling (e.g., searching specific function by keyword, listing functions that call a certain function) operations to dynamically explore repositories.
However, the latter approach places higher demands on LLMs: beyond multi-step
reasoning, they require sophisticated tool-calling abilities to explore code effectively rather than walking aimlessly through the repository.

Agentic codegen systems such as OpenHands –\citet{wang2025openhandsopenplatformai} utilize agentic LLMs that can iteratively read code, call tools, and update their hypotheses – abilities that are well suited to fault localization.
It requires an agent (tool calling LLM) to write legit bash commands that will be executed in the repository dir.
However often times an agent pipes multiple commands such as "ls", "find", "cat" and more together to get a representative information, which is not the case for human developers when localizing a bug.
Moreover these systems often times spend half of the overall interaction turns on localization task, essentially occupying unnecessary context for the patch generator component.

Given these shortcomings there is a need to have a specialized localization system that will work separately from the other components by simply providing the locations of the edit, leaving room for the other 2 components to actively iterate and refine the solution.
\textbf{LocAgent} – emulates agentic behavior to produce multi-turn localization trajectories at scale with open-weight models and a finite tool set, then distill a smaller student model to pinpoint buggy lines.





% \textbf{Data generation for distillation.} We use Qwen3-32B as the generator and leverage its tool-calling to read files (via a simple file-reader, e.g., \texttt{function\_calling}). Inputs are PR problem statements and repository structures. Over $N$ turns, the LLM receives the problem statement, a short “conversational memory” summary it chooses to retain, and (when applicable) outputs from the previous tool call. At each turn it produces a reasoning trace and optional tool call; upon identifying the buggy lines, it emits \texttt{<answer>...</answer>} and the conversation terminates. A final verification step checks whether the predicted lines are correct.

% \textbf{Learning and evaluation.} We distill the collected trajectories into a student model tailored to fault localization and evaluate on \textsc{SWE-bench-verified}, \textsc{SWE-bench-live}, and \textsc{Loc-Bench}.


\section{Related Work}

OpenHands \citep{wang2025openhandsopenplatformai} simulates the repo exploration in a sandbox, it uses tools such as "open\_file", "goto\_line", "scroll\_down" that actually open the files and wait for an action in the next turn.

SWE-Agent \citep{yang2024sweagentagentcomputerinterfacesenable} hsa specialized tools, including an interactive file viewer, search functionalities, and edit tools for the open file, in additional to the standard Linux Bash commands, similar to OpenHands\citep{wang2025openhandsopenplatformai}.

Tool integrated reinforcement learning repo \cite{ma2025toolintegratedreinforcementlearningrepo}.

PathPilot \cite{li2025patchpilotcostefficientsoftwareengineering}.


\section{Method}

We present LocAgent, a multi-turn interaction system designed for precise issue localization in software repositories. Unlike existing approaches that rely on unbounded bash commands or complex tool sets, LocAgent employs a focused set of four tools within a carefully managed dialogue framework to efficiently identify code locations requiring modification.

\subsection{Multi-Turn Interaction Framework}

LocAgent models issue localization as a sequential decision-making process where an autonomous agent iteratively explores the codebase through structured tool interactions. The agent operates within a dialogue framework, maintaining conversation history that captures both its reasoning process and information gathered from the repository.

The agent follows a perception-reasoning-action cycle: (1) receiving the current state comprising the problem description and accumulated tool outputs, (2) analyzing available information using a large language model to determine the next exploration action, and (3) executing one of four specialized tools to gather additional repository information.

\subsection{Tool Design}

We provide four carefully designed tools that enable comprehensive yet efficient repository exploration:

\textbf{File Viewing} retrieves source code content with optional line range specification, enabling focused examination of specific code sections while managing output size through configurable limits.

\textbf{Repository Structure} generates a hierarchical view of the codebase, intelligently filtered to show only relevant source files while excluding non-essential directories such as tests, build artifacts, and documentation.

\textbf{Code Search} performs case-insensitive keyword search across the repository, returning contextual snippets (20 lines before/after) around matches to facilitate understanding of code usage patterns. Results are ranked by match frequency and limited to maintain manageable output sizes.

\textbf{Dependency Analysis} constructs import relationship graphs to understand module interconnections, helping identify related components that may require modification.

\subsection{Context Management}

A critical challenge in multi-turn systems is managing growing dialogue history within language model context limitations. We developed a comprehensive context management strategy with multiple components:

\subsubsection{Token Estimation}
We implement context-aware token estimation that differentiates between code and natural language. Code segments require approximately 3.5 characters per token due to operators and formatting, while natural language averages 4.5 characters per token. This differentiated estimation enables accurate context planning.

\subsubsection{Truncation Strategies}
When dialogue history approaches context limits, we employ one of several strategies:

\textbf{Sequential Truncation} progressively removes oldest interaction turns while preserving recent context, suitable for problems requiring detailed tracking of recent explorations.

\textbf{Bookend Truncation} (default) retains the initial problem statement and most recent interactions while removing intermediate turns, ensuring the agent maintains problem awareness while focusing on current exploration state.

\textbf{Adaptive Truncation} dynamically selects between strategies based on conversation characteristics, with multi-level fallbacks for extreme cases requiring aggressive reduction.

\subsubsection{Proactive Monitoring}
Rather than reacting to context overflow errors, our system proactively monitors token usage before each interaction, calculating expected consumption and applying appropriate truncation strategies preemptively. A configurable safety margin (default: 90\% of maximum context) ensures reliable operation.

\subsection{Loop Detection and Intervention}

Multi-turn systems risk entering repetitive behavior patterns. We address this through comprehensive loop detection and intervention:

The system continuously monitors tool usage patterns, maintaining a sliding window of recent actions. When consecutive identical tool calls exceed a threshold (default: 3), it identifies a behavioral loop and injects carefully crafted intervention messages suggesting alternative exploration strategies.

Beyond immediate intervention, the system logs repetitive patterns to identify common failure modes, informing both prompt engineering improvements and architectural enhancements.

\subsection{Repository Filtering}

Efficient exploration requires intelligent filtering to focus on relevant code:

\textbf{File Type Selection}: We restrict exploration to source code files most likely to contain implementation logic (Python and configuration files for Python repositories).

\textbf{Directory Filtering}: We exclude directories unlikely to contain production code including testing infrastructure, build artifacts, documentation, and third-party dependencies.

\textbf{Adaptive Preservation}: Based on empirical analysis, we preserve commonly excluded directories that frequently contain relevant code: utility modules, library directories, migration files, and configuration directories.

\subsection{Final Turn Optimization}

To maximize success rates given finite interaction budgets (default: 20 turns), we implement final turn optimization. When approaching the maximum allowed interactions, the system injects specialized instructions requiring the agent to synthesize accumulated information into concrete location predictions. This transforms potential failures due to interaction limits into partial successes based on accumulated evidence.

\subsection{Robustness Features}

Real-world deployment requires handling various failure modes gracefully:

\textbf{Format Flexibility}: The system accepts multiple tool invocation formats, including properly structured commands, partial specifications, and implicit requests, accommodating variations in language model outputs.

\textbf{Progressive Degradation}: When encountering errors, the system attempts multiple recovery strategies: suggesting similar files when exact matches aren't found, providing informative error messages with actionable suggestions, and maintaining operation despite individual tool failures.

\textbf{Implicit Tool Detection}: When agents fail to use proper tool call syntax, the system employs fallback strategies including JSON detection after reasoning blocks, simple file path extraction, and search query pattern matching.

\subsection{Evaluation Framework}

We implement comprehensive evaluation mechanisms including ground truth extraction from repository patches, multi-level metrics (file-level accuracy, line-level precision, coverage analysis), and systematic failure categorization to identify areas for improvement.

The modular design enables controlled experimentation while maintaining robustness across diverse repository structures and problem types, with all enhancements maintaining backward compatibility for ablation studies.


\bibliography{custom}

\appendix

\section{Example Appendix}
\label{sec:appendix}

This is an appendix.

\end{document}
